\section{B. Adelic approximations, completions, and exactness}
\subsection{The definition actually given in \S52}

\begin{statement}{B.1}{Proved}{The axioms of a Galois--Teichm\"uller structure}
The final provisional definition in \MS{52}{750--752}{205--209}
consists of inclusions
\[
 \Pi_0\subset\mathcal G_0,\quad \Pi\subset\mathcal G\subset\mathcal M,
 \qquad \Pi=\widehat{\Pi_0},\quad \mathcal G=\overline{\mathcal G_0},
\]
with the second row profinite, the first discrete, $\Pi,\mathcal G$ normal
in $\mathcal M$, the specified peripheral $L$-structure invariant under
$\mathcal M$, and $\mathcal G_0$ normalizing $\Pi_0$. Its cyclotomic
character factors through $\Gamma=\mathcal M/\mathcal G$ and has open
image in $\widehat\Z^\times$. We do not require
$\widehat{\mathcal G_0}\cong\mathcal G$, which the manuscript expressly
does not require at scan206.

A Lie approximation is an affine finite-type $\Q$-group $G$, a character
$\chi_G:G\to\Gm$, and a continuous homomorphism
$\rho:\mathcal M\to G(\mathbb A_f)$ such that
\[
 \chi_G\rho=\chi,\qquad \rho(\mathcal G_0)\subset G(\Q),
\]
where $\Q$ is embedded diagonally, and such that no proper closed
$\Q$-subgroup contains $\rho(\mathcal M)$ adelically.
These are conditions $1^\circ,2^\circ,3^\circ$ forte, and $4^\circ$
of scans202--208.
\end{statement}
\begin{proof}
Diagram (35), visually checked at scan208, factors through $G(\Q)$;
the ambient representation has target $G(\mathcal A_\Q)$, not
$G(\Q_\ell)$ for a single prime. The strengthening from $\Pi_0$
to $\mathcal G_0$ is explicit in scans204--208.
The arithmetic fundamental-group sequence of a hyperbolic curve over a
finitely generated field of characteristic zero supplies the case
$\mathcal G=\Pi$: choose a complex comparison and its discrete
peripheral group. For the universal smooth curve, comparison with
the complex topological fibre and moduli stack supplies the
geometric groups and their closures. The étale homotopy sequence
gives normality; conjugation on a puncture's inertia is cyclotomic,
because an $n$th-root local parameter is multiplied by an $n$th
root of unity. The character is trivial geometrically. The algebraic
constants in the base field have finite degree over $\Q$, hence
their intersection with $\Q(\mu_\infty)$ is finite; the cyclotomic
image is open. These checks do not assert the congruence-subgroup
property for mapping class groups. With no cusps a cyclotomic
character is extra data, not recoverable from an empty list.
\end{proof}
\begin{statement}{B.2}{Proved}{A universal pro-object of admissible targets}
For fixed B.1 data, let morphisms of admissible targets be
$\Q$-group morphisms compatible with $\rho$ and $\chi_G$.
This category is cofiltered and has at most one morphism between
two objects. Its inverse limit $\mathscr E$ is a pro-affine
$\Q$-group universal for the specified admissible finite-type
targets. This is not an identification with the absolute
proalgebraic envelope of $\Pi_0$ or its Malcev completion.
\end{statement}
\begin{proof}
The object $\Gm$ with $\rho=\chi$ exists. Given two targets take
the scheme-theoretic $\Q$-Zariski closure of the image in their
product over $\Gm$. Intersecting the closed subgroups containing
the image gives a closed subgroup; finite type follows from
noetherianity of the ambient coordinate ring. This is a common
refinement. Continuity, rationality on $\mathcal G_0$, and the
character condition are preserved, and minimality is built into
the closure. Two morphisms agreeing on the adelic image agree on
its Zariski closure and hence everywhere. Choose a set of
representatives of the finite-type coordinate algebras and
representations. The direct limit of their coordinate Hopf
algebras represents the claimed inverse limit. This proves its
universal property without imposing unipotence or weights.
\end{proof}

\begin{statement}{B.3}{Refuted}{Lead 1 as an identification with a nilpotent tower}
B.1 is not the replacement of the geometric groups by nilpotent
or solvable completions. It specifies neither a first nor a second
stage. A pair of first two approximations therefore requires a
chosen filtration or subcategory of targets.
\end{statement}
\begin{proof}
For the elliptic modular arithmetic fundamental group take the
adelic Tate-module representation in $\GL_{2,\Q}$, with determinant
as character. Topological comparison sends the geometric discrete
group $\SL_2(\Z)$ to its usual rational matrices. Its Zariski
closure is $\SL_{2,\Q}$: it contains both upper and lower
one-parameter unipotent subgroups, whose brackets generate
$\mathfrak{sl}_2$. The determinant has open cyclotomic image,
so the full closure is $\GL_{2,\Q}$. This is an admissible target
with non-solvable connected geometric group. A construction
making the whole geometric group prounipotent cannot have this
image. Weighted completion, relative completion with a fixed
reductive representation, and genus-zero or elliptic $\GT$
specify different universal problems; the label Lie approximation
alone supplies no comparison map. The strongest corrected
interpretation is the category of B.1 and its pro-object B.2.
\end{proof}

\begin{statement}{B.4}{Proved}{Two explicitly chosen Malcev stages}
For types $(0,3),(0,4),(1,1)$ the geometric discrete group is
free of rank $r=2,3,2$. Choose rational Malcev nilpotence-class
truncation. Its first stages are
\[
 U_1=\Ga^r,\qquad
 U_2=\bigl(V\oplus\Lambda^2V,*\bigr),\qquad
 (v,a)*(w,b)=(v+w,a+b+\tfrac12v\wedge w).
\]
The Lie bracket is $[v,w]=v\wedge w$ and $\Lambda^2V$ is central;
lower-central degrees are $1,2$. For punctured spheres the graded
Galois representations are $\Q_\ell(1)^r$ and
$\Q_\ell(2)^{\binom r2}$; for a punctured elliptic curve they are
$T_\ell E\otimes\Q_\ell$ and $\Q_\ell(1)$. Permutations of punctures
act on $V$ by permuting peripheral classes subject to their sum
being zero, and act by the exterior square in degree two.
\end{statement}
\begin{proof}
The displayed Lie algebra has the universal property of the
free class-two Lie algebra: arbitrary images of a basis of $V$
extend uniquely, and all brackets of length three vanish.
Truncated Baker--Campbell--Hausdorff gives the group law;
bilinearity verifies associativity and inverses are $(-v,-a)$.
Cutting the three surfaces along standard arcs gives their free
fundamental groups. Functoriality determines the graded action
from abelianization. Peripheral inertia is $\Z_\ell(1)$; for the
torus the Weil pairing identifies $\Lambda^2T_\ell E$ with
$\Z_\ell(1)$. This determines the graded action, not a canonical
splitting of a filtered genuine action. Base-point cocycles and
integral models are additional choices (Parts D, E, G); in
particular the factor $1/2$ requires attention at $2$.
\end{proof}

\subsection{A covering-exact family within the scope of \S50}

\begin{statement}{B.5}{Proved}{Virtually abelian quotients and finite covers}
For a discrete group $F$ put
\[
 V(F)=\varprojlim_{K\triangleleft F,\ [F:K]<\infty}F/[K,K].
\]
Equip this inverse limit with its pro-discrete topology.
If $F\twoheadrightarrow H$ is a finite quotient with kernel $N$,
then $1\to V(N)\to V(F)\to H\to1$ is exact as an inverse limit
of discrete groups. For finitely generated free $F$ the \S50 envelopes apply to
every stage, and their inverse system is also exact for this
finite covering.
\end{statement}
\begin{proof}
Restrict to $K\subset N$, a cofinal subsystem. It is also cofinal
among the finite-index normal subgroups of $N$: intersect finitely
many $F$-conjugates of such a subgroup. At each stage,
$1\to N/[K,K]\to F/[K,K]\to H\to1$ is exact. Kernels commute
with inverse limits, and each element of $H$ has compatible lifts
given by a fixed representative in $F$.

If $F$ is free, $K^{\ab}$ is a finite-rank lattice. The two groups
$N/[K,K]\subset F/[K,K]$ have this same normal lattice.
In the explicit \S50 construction, each finite scheme stage is
a disjoint union of affine spaces with multiplication determined
by an integral extension cocycle. The subgroup over $N/K$ is
the union of the corresponding components, hence its envelope
is exactly the fibre over the identity of the constant group
$H$. Transition maps preserve these component decompositions;
inverse limits prove the assertion for the envelopes. Without
the lattice hypothesis only the group inverse-limit statement
has been asserted here.
\end{proof}

\begin{statement}{B.6}{Proved}{Fundamental localizers of the specified completions}
For abelianization, the metabelian quotient, a fixed positive
nilpotence-class quotient, pro-$\ell$ completion, rational
prounipotent completion, and the virtually abelian pro-object
$V$, the class $W_{\mathfrak a}$ in the
prompt satisfies two-out-of-three, retract closure, the
terminal-object axiom, and relative Theorem A. The literal
\S52 construction requires richer arithmetic data and does not
define a functor of an arbitrary bare fundamental group.
\end{statement}
\begin{proof}
Use all coefficient groups $H$ which are, respectively, abelian,
metabelian, nilpotent of the fixed class, finite $\ell$-groups,
groups $U(\Q)$ for finite-dimensional unipotent $\Q$-groups,
or all virtually abelian discrete groups for $V$.
For a connected space $X$, $\operatorname{Map}(X,BH)$ is the
nerve of the groupoid of homomorphisms $\pi_1X\to H$ and their
conjugations. The completion's universal property identifies
these representations, and also their centralizers: finite
images are unchanged; in the unipotent case the Zariski closures
of the images are unchanged. Conversely equivalence on all
these mapping spaces determines the quotient or pro-object by
its universal property: the mapping-space comparison commutes
with evaluation at the base point to $BH$, and its homotopy
fibre is the discrete set of based homomorphisms. For $V$,
a map to a discrete virtually abelian group factors through
$F/[K,K]$ for the inverse image $K$ of an abelian subgroup
of finite index; all maps out of $V$ here are continuous.
Including discrete sets as targets
detects $\pi_0$. Thus $W_{\mathfrak a}$ is exactly the class
inverted by all these mapping-space functors.

Two-out-of-three and retract closure follow. Mapping out of a
homotopy colimit gives a homotopy limit. For a triangle
$A\to B\to S$ of small categories, the Grothendieck construction
gives
\[
 BA\simeq\operatorname*{hocolim}_{s\in S}B(A/s),\qquad
 BB\simeq\operatorname*{hocolim}_{s\in S}B(B/s).
\]
The first comparison follows by projecting the category of
pairs $(a,f(a)\to s)$ to $A$: the arrows out of $f(a)$ form
a category with initial object its identity. The second is
identical. An objectwise equivalence after mapping to every
test target therefore gives the required global equivalence.
This proves relative Theorem A. A terminal object contracts the
nerve, proving its axiom. This argument does not imply that all
coverings are intrinsic fibrations for these localizers.
\end{proof}

\begin{statement}{B.7}{Proved}{Typed comparisons for finite quotients}
Let $F=F_2$, $F\twoheadrightarrow H$, $d=|H|$, and $N$ its kernel.
The prompt's comparison map exists only when this finite quotient
factors through the chosen approximation. In that case it is
surjective, with the following images and kernels:
\[
\begin{array}{c|c|c}
\text{approximation}&\text{image}&\text{kernel}\\ \hline
(-)^{\ab}&N/F'&(N\cap F')/N'\\
(-)/(-)''&N/F''&F''/N''\\
(-)/\gamma_{c+1}&N/\gamma_{c+1}F&
                    \gamma_{c+1}F/\gamma_{c+1}N .
\end{array}
\]
The required conditions are respectively $H$ abelian, $H''=1$,
and $\gamma_{c+1}H=1$. The first kernel is free abelian of rank
$d-1$. None of these three maps is an isomorphism for $d>1$.
Pro-$\ell$ completion is exact for all finite $\ell$-group
quotients; other finite groups cannot be its continuous quotients.
A connected prounipotent $\Q$-group has no nontrivial finite
étale quotient.
\end{statement}
\begin{proof}
The quotient formulas are the third isomorphism theorem.
Schreier gives $\rank N=d+1$. In the abelian row its image is
a rank-two lattice. In the nilpotent row the target has
rational abelianization of rank two, since its Malcev Lie algebra
is the free class-$c$ Lie algebra on two generators, whereas the
source has abelianization of rank $d+1$.

For the metabelian row choose $v\in N\cap F'$ nonzero in
$N^{\ab}$, using the same rank inequality. The Magnus matrix
embedding identifies the derived subgroup of a free metabelian
group with a submodule of a free module over the integral domain
$\Z[N^{\ab}]$. Conjugation by $v$ is multiplication by a
nonidentity monomial, so $v$ fails to centralize some
$w\in N'/N''$. Thus $[v,w]\ne1$ modulo $N''$, but
$v,w\in F'$ imply $[v,w]\in F''$. This proves a nonzero kernel.
Explicitly the Magnus embedding sends a free basis $t_i$ to
$\left(\begin{smallmatrix}T_i&e_i\\0&1\end{smallmatrix}\right)$
over the Laurent-polynomial module and has kernel $N''$.

For an $\ell$-group $H$ the induced pro-$\ell$ topology on $N$
is its full pro-$\ell$ topology: intersect the finitely many
$F$-conjugates of a normal subgroup of $\ell$-power index in
$N$. Its quotient and the quotient $F/N$ have $\ell$-power
orders, so the intersection is open in the pro-$\ell$ topology
of $F$. Completion is therefore exact. The converses follow
from the elementary finite-quotient properties stated. For
unipotent groups, connectedness forces every map to a finite
constant group to be trivial at each finite-dimensional stage.
The missing absolute-unipotent map is not an inexact map.
\end{proof}

\begin{statement}{B.8}{Proved}{The Fermat cover and relative completion}
For $N=\ker(F_2\to(\Z/n)^2)$, $n>1$, the absolute map
$N^{\un}_\Q\to (F_2^{\un})_\Q$ is surjective. On abelianizations
its source, target, and kernel have dimensions $n^2+1$, $2$,
and $n^2-1$. Relative completion with respect to
$H=(\Z/n)^2$ instead gives
\[
 1\to N^{\un}_\Q\to\mathscr R_H(F_2)\to H_\Q\to1.
\]
Its abelian unipotent kernel has dimension $n^2+1$, and
$H_1(N,\Q)\cong\Q\oplus\Q[H]$ as an $H$-module.
\end{statement}
\begin{proof}
Schreier gives $\rank N=1+n^2$. A finite-index subgroup has
Zariski-dense image in every connected unipotent quotient of
$F_2$: finitely many cosets of its image closure cover the group,
so the closure has finite index and equals the connected group.
Equivalently the map on rational $H_1$ is onto, and the free
pronilpotent Lie algebra universal property gives surjectivity.

Construct the relative group as the inverse limit of unipotent
extensions of $H_\Q$ with a lift of $F_2\to H$. Its kernel
has exactly the universal property of $N^{\un}_\Q$: induction
over the finitely many cosets embeds every unipotent
$N$-representation in a finite-component representation of
$F_2$, and restriction gives the inverse comparison of
universal properties. The cellular chains of the regular
covering of a two-petalled rose give
$0\to H_1\to\Q[H]^2\to\Q[H]\to\Q\to0$.
Semisimplicity of $\Q[H]$ proves the module formula.
Its nontrivial finite-order eigenspaces cannot occur in a
unipotent action, because a finite-order unipotent matrix in
characteristic zero is the identity. This is the precise
information forgotten by absolute completion; a group of
conjugator parameters of dimension six is a different object.
\end{proof}

\begin{statement}{B.9}{Proved}{Intrinsic finite coverings}
Finite coverings intrinsic for $W_{\mathrm{pro}\text{-}\ell}$
are precisely those with $\ell$-group monodromy. Those intrinsic
for rational unipotent completion are precisely those with
trivial monodromy. A finite covering of $BF_2$ with nontrivial
monodromy is intrinsic for neither abelianization, the metabelian
quotient, nor any fixed positive nilpotence-class approximation.
Every finite covering is intrinsic for the virtually abelian
pro-quotient functor $V$ of B.5.
\end{statement}
\begin{proof}
Pro-$\ell$ equivalences preserve finite $\ell$-sets and the
completions of their stabilizers by the induced-topology
argument of B.7. This proves preservation after every base
change. If finite monodromy is not an $\ell$-group, an element
of prime order $q\ne\ell$ acts nontrivially on the fibre.
Pull back along a loop representing it. Multiplication by $q$
on $B\Z$ is a pro-$\ell$ equivalence but splits a nontrivial
orbit into singletons, failing the component test. For rational
unipotent completion the same proof uses the order of any
nontrivial monodromy element.

For the abelian and nilpotent assertions pull back to a two-petalled rose by
sending $x$ to a nontrivial monodromy loop and $y$ to the
identity. One orbit gives a cyclic $m$-sheeted cover, $m>1$.
The endomorphism $\phi(x)=x$, $\phi(y)=yx^{-1}y^{-1}xy$ is the
identity on abelianization and an automorphism on every
nilpotent quotient: it is the identity on its associated
graded, so induction along the central extensions applies.
It preserves cyclic monodromy. A Schreier basis is
$a=x^m$, $b_i=x^iyx^{-i}$, $0\le i<m$. On its abelianization,
$a$ is fixed and $b_i\mapsto2b_i-b_{i-1}$, with cyclic
subscripts. Here the group commutator used in the prompt is $[x,y]=x^{-1}y^{-1}xy$.
The determinant is $2^m-1\ne\pm1$. Thus the pullback is
not an abelian equivalence and not a fixed-class nilpotent
equivalence.

For the metabelian assertion pull back instead along $F(x,y,z)$,
with $x$ carrying the same cyclic monodromy and $y,z$ trivial.
Set $w=[[x,y],[x,z]]$ using $[a,b]=a^{-1}b^{-1}ab$, and let
$\psi$ fix $x,z$ and send $y$ to $yw$. It is the identity on
$F/F''$, and is compatible with this base change. In the
cyclic cover choose Schreier generators
$a=x^m,y_i=x^iyx^{-i},z_i=x^izx^{-i}$.
On its abelianization $\psi$ is the identity. In its abelianized
Fox Jacobian specialize the variables of $a,y_i$ to $1$ and
write the variables of $z_i$ as $Z_i$.
The nontrivial block has row $i$
\[
 R_i e_i-(R_i-1)e_{i-1},\qquad R_i=Z_{i-1}/Z_i.
\]
Indeed $[x,y]$ at vertex $i$ rewrites as
$y_{i-1}^{-1}y_i$, and the Fox derivative of its commutator
with $z_{i-1}^{-1}z_i$, after this specialization, is
$(R_i-1)(e_i-e_{i-1})$. Thus the determinant is
$1-\prod_i(R_i-1)$.
Specializing $Z_i=t^i$ gives
\[
 1-(t^{m-1}-1)(t^{-1}-1)^{m-1}.
\]
Its highest and lowest exponents are $m-1$ and $-(m-1)$,
with nonzero coefficients. It is not a Laurent monomial.
An automorphism of the free metabelian group would have an
invertible Fox Jacobian by the chain rule, hence a unit
determinant. This contradiction proves failure on the cover.

Finally $V$ is characterized by homomorphisms into virtually
abelian groups. A representation of a finite cover into such
a group $A$ induces a representation downstairs into
$A^\Omega\rtimes\operatorname{Sym}(\Omega)$, with prescribed
permutation action on the finite fibre $\Omega$. This group is
virtually abelian. The induction correspondence, including
conjugations, identifies the representation groupoids upstairs
with the fibre of this downstairs representation groupoid over
the monodromy. A $V$-equivalence preserves both groupoids and
this fibre. Applying the same argument componentwise proves
that the covering is intrinsic. For arbitrary category bases,
a fundamental-group loop is represented by a finite zigzag
in the nerve. Its subdivision is a small category and maps
to the original category, so the same loop tests are valid
base changes; a literal functor $B\Z\to B$ is not required.
These arguments separate
exactness on one cover from universal preservation.
\end{proof}

\subsection{Special derivations and integral lifting}

\begin{statement}{B.10}{Proved}{Exact modules and the three tripod dimension sequences}
Let $L(V)$ be the free Lie algebra over $\Z$, with
$V=\Z^q/\Z(1,\ldots,1)$, $r=q-1$, and let $\mathfrak s_{r,n}$
denote degree-$n$ tangential special derivations modulo inner
derivations. For $n\ge2$ there is an $S_q$-equivariant exact sequence
\[
 0\longrightarrow\mathfrak s_{r,n}
 \longrightarrow V\otimes L_n(V)
 \xrightarrow{[\, ,\, ]}L_{n+1}(V)\longrightarrow0.
\]
In degree one one additionally divides by the $q$ peripheral
centralizer lines. Consequently
\[
 \dim\mathfrak s_{r,n}=r\ell_n(r)-\ell_{n+1}(r)
                    -(r+1)\delta_{n,1},
 \qquad
 \ell_n(r)=\frac1n\sum_{d\mid n}\mu(d)r^{n/d}.
\]
For the tripod, the dimensions through weight fourteen are:
\begin{center}\begin{tabular}{r|rrrr}
\toprule Degree&All special outer&Raw cycles&Special invariants&Odd free Lie\\\midrule
1 & 0 & 1 & 0 & 0\\
2 & 0 & 0 & 0 & 0\\
3 & 1 & 1 & 1 & 1\\
4 & 0 & 1 & 0 & 0\\
5 & 3 & 2 & 1 & 1\\
6 & 0 & 3 & 0 & 0\\
7 & 6 & 6 & 2 & 1\\
8 & 4 & 10 & 1 & 1\\
9 & 13 & 19 & 4 & 1\\
10 & 12 & 33 & 2 & 1\\
11 & 37 & 62 & 9 & 2\\
12 & 40 & 112 & 7 & 2\\
13 & 99 & 210 & 21 & 3\\
14 & 140 & 387 & 24 & 3\\
\bottomrule\end{tabular}\end{center}
Here raw cycles means
$\psi(x,y)+\psi(y,x)=0$ and
$\psi(x,y)+\psi(y,z)+\psi(z,x)=0$, $z=-x-y$.
The special-invariant column instead imposes antisymmetry and
\[
 J(\psi)=[x,\psi(x,y)]-[x+y,\psi(-x-y,y)]=0.
\]
The last column is the free graded Lie algebra on one generator
in each odd degree at least three; no equality with all of
$\grt_1$ is built into this computation.
\end{statement}
\begin{proof}
Write $D(x_i)=[a_i,x_i]$ and
$D(\sum x_i)=[a_0,\sum x_i]$. Subtracting the inner derivation
defined by $a_0$ makes the sum zero. Brackets with the free
generators span $L_{n+1}$ integrally: repeatedly applying Jacobi
moves one generator to the first position of a bracket.
The centralizer of a free generator in a free Lie algebra is its
one-dimensional span. This follows, for example, by embedding in
the free associative algebra and comparing the first letter of
each word in $[x_i,a]$. Thus the $a_i$ are unique for $n\ge2$,
and the residual inner ambiguity centralizing the sum is also
zero in those degrees. Keeping all $q$ peripheral conjugators
and then dividing by their common inner change identifies the
coefficient module equivariantly with
$(\Z^q/\Z)\otimes L_n(V)$. This proves the sequence; in degree
one the $q$ redundant peripheral conjugators give the stated
subtraction. Integral surjectivity also proves that reduction
modulo a prime computes the same kernel.

The multigraded Poincaré--Birkhoff--Witt identity, followed by
Möbius inversion, gives, for every permutation $g$,
\[
 \operatorname{tr}(g\mid L_nV)
 =\frac1n\sum_{d\mid n}\mu(d)
       \operatorname{tr}(g^d\mid V)^{n/d}.
\]
Here $\operatorname{tr}(g^d\mid V)$ is the number of fixed
punctures of $g^d$ minus one. Subtract the two terms in the
exact sequence and average over $S_q$ to compute invariants.
For raw tripod cycles, the antisymmetric vector with zero
three-cycle norm is one line in each standard $S_3$-summand;
its multiplicity is
$(\operatorname{tr}(1\mid L_n)-\operatorname{tr}((123)\mid L_n))/3$.

For clarity on the extra equation, normalize an outer
derivation by making its $x$-image zero and writing its
$y$-image as $[y,\psi]$. Commutation, modulo inner
derivations, with the transposition of $x,y$ gives
antisymmetry. Commutation with the cyclic permutation of
$x,y,-x-y$, after subtracting the unique inner discrepancy,
is exactly $J(\psi)=0$. These are derivation equations; simply
linearizing two word products and forgetting the conjugation
action would omit this second condition.

The programs construct standard Lyndon brackets in the
associative tensor algebra. They form the antisymmetry,
cycle, and $J$ matrices, compute their kernels by exact
rational row reduction through weight thirteen, and compare
with the independent character formula. Weight fourteen uses
the exact character formula after a direct expanded matrix
attempt exhausted memory; this is not a rank guessed modulo
a prime. The odd-generator column is independently obtained
by equating coefficients in
\[
 \prod_{n\ge1}(1-t^n)^{-d_n}
       =\frac1{1-t^3-t^5-t^7-\cdots}.
\]
All coefficients, rational kernels through thirteen, and the
failed direct attempt are retained in the computational output.
\end{proof}

\begin{statement}{B.11}{Proved}{Every integral special leading term lifts}
Let $F$ be free pro-$\ell$ on $x_1,\ldots,x_r$ and let
$x_0=(x_1\cdots x_r)^{-1}$. Every homogeneous special
derivation of $\gr_\gamma F=L(\Z_\ell^r)$ is the leading
term of a continuous automorphism conjugating each $x_i$,
including $x_0$, to itself. The assertion holds for every
prime, including $2$ and $3$, and passes to outer classes.
\end{statement}
\begin{proof}
For degree $n$, choose lifts $A_i\in\gamma_nF$ of the
conjugators $a_i$. They prescribe tentative images
$A_i x_i A_i^{-1}$. The special identity says that
\[
 \left(\prod_{i=1}^r A_i x_i A_i^{-1}\right)
      \left(A_0(x_1\cdots x_r)A_0^{-1}\right)^{-1}
       \in\gamma_{n+2}F .
\]
Suppose the error is in $\gamma_mF$, $m\ge n+2$.
Multiplying the $A_i$ by elements of $\gamma_{m-1}F$
changes its class in $\gamma_m/\gamma_{m+1}$ by
\[
 (b_0,b_1,\ldots,b_r)\longmapsto
      \sum_i[b_i,x_i]-[b_0,\sum_i x_i].
\]
This map is onto over $\Z_\ell$, already with $b_0=0$,
by the integral bracket-surjectivity proved in B.10.
We can therefore cancel the error without changing the
prescribed leading term. Iterating produces compatible
conjugators and an exact product identity.

The iteration converges in the actual pro-$\ell$ topology:
every finite $\ell$-group quotient has finite nilpotence
class, so eventually all corrections vanish in that quotient.
The images induce the identity on the Frattini quotient,
hence generate $F$ by the pro-$\ell$ Burnside basis theorem.
The resulting continuous endomorphism is surjective and
therefore injective: a finitely generated profinite group
has only finitely many subgroups of each index, and a
surjective endomorphism permutes them, which forces its
kernel into every open normal subgroup. This proves that it
is an automorphism. Inner conjugations realize the inner
derivations, so quotienting gives the outer assertion.
No exponential map or denominator estimate is used.
\end{proof}

\begin{statement}{B.12}{Proved}{Coprime symmetry lifting and strict pro-$\ell$ inclusion}
For $\ell\ge5$, every $S_3$-invariant integral special outer
leading term for the tripod lifts to the centralizer of the
geometric $S_3$ in the special outer automorphism group.
With multiplier one this centralizer is the pro-$\ell$
group of solutions of the two Drinfeld cycle equations.
It strictly contains the subgroup also satisfying the
pentagon.
\end{statement}
\begin{proof}
Let $P_n$ be the subgroup of special outer automorphisms
whose first possible derivation degree is $n$. It is
pro-$\ell$: in each finite characteristic $\ell$-group
quotient, the kernel of the action on the Frattini quotient
is an $\ell$-group; taking inverse limits proves the claim.
B.11 identifies $P_n/P_{n+1}$ with the integral module
in B.10. If $a$ has invariant image, the discrepancies
$a^{-1}g(a)$ form a nonabelian $1$-cocycle of $S_3$ in
$P_{n+1}$. Its cohomology class is trivial. To see this
without an averaging assumption on a nonabelian group,
work in every finite $S_3$-stable quotient of $P_{n+1}$:
successively average over the central elementary-abelian
layers of its invariant central series. Multiplication
by $6$ is invertible on each layer, so all the additive
$H^1$ obstructions vanish. Equivalently the complements
to the finite normal $\ell$-group are conjugate.
The nonempty solution sets at finite level are compact
and have the finite-intersection property. Their inverse
limit supplies a single conjugating correction. Thus
$P_n^{S_3}\to(P_n/P_{n+1})^{S_3}$ is onto.

Here is the precise group-equation identification. An
outer special automorphism has a unique normalized
representative
\[
 x\longmapsto x,\qquad y\longmapsto f(x,y)^{-1}y f(x,y),
 \qquad f\in F',
\]
after removing the two abelian exponents of its conjugator.
Uniqueness follows from the procyclic centralizers of
$x$ and $y$ in a free pro-$\ell$ group. The word
calculations in \citet[Propositions 3, 4 and Theorem 2]{HarbaterSchneps2000}
then give
\[
 f(y,x)f(x,y)=1,\qquad
 f(z,x)f(y,z)f(x,y)=1,\quad xyz=1.
\]
More explicitly, the first equation is precisely that
the commutator with the transposition is inner; after it
is imposed, the commutator with the three-cycle is inner
precisely when the second product is central in the
peripheral subgroup. Its abelianization is zero, so that
central product is the identity. Conversely the same
computations exhibit the required inner commutators.
The proof uses only free-group substitution, abelianization,
and the stated peripheral centralizers; all remain valid
in the free pro-$\ell$ group. The finite symmetry is
retained through the completion relative to $S_3$.
Taking the absolute pro-$\ell$ completion of $B_3$ and
discarding that symmetry would not justify this argument.

Let $b_w$ denote the standard Lyndon bracket for $w$.
The explicit primitive integral vector
\[
 \begin{split}
 q_7={}&3b_{xxxyxyy}+3b_{xxxyyxy}+b_{xxyxyxy}\\
       &-3b_{xxyxyyy}-3b_{xxyyxyy}
        -3b_{xxyyyxy}-b_{xyxyxyy}
 \end{split}
\]
satisfies antisymmetry and $J(q_7)=0$, as direct integer
expansion verifies. In the infinitesimal pure braid
representation on $L(a,b,c)$, its pentagon defect applied
to $a$ has coefficient $16$ on $a^5bcb$.
The program first verifies all defining braid relations,
then expands all five pentagon terms; the full output
has $2570$ nonzero coefficients and is attached.
Thus its defect is nonzero over every $\Q_\ell$, and is
already nonzero modulo every $\ell\ge5$.
Lift $q_7$ by the preceding symmetry argument.
Higher-degree corrections cannot change the first
nonzero pentagon degree. The lifted element therefore
satisfies the two cycles and fails the pentagon.
\end{proof}

\begin{statement}{B.13}{Proved}{The modular cohomology obstruction table}
Let $M_n=\mathfrak s_{2,n}\otimes\F_p$ with its integral
geometric $S_3$-action reduced modulo $p$. In degrees
$1,\ldots,9$, both $H^1(S_3,M_n)$ and $H^2(S_3,M_n)$ have
the following dimensions:
\[
\begin{array}{c|rrrrrrrrr}
n&1&2&3&4&5&6&7&8&9\\ \hline
p=2&0&0&1&0&1&0&2&0&3\\
p=3&0&0&0&0&0&0&0&1&0
\end{array}
\]
These are obstruction \emph{spaces}, not evaluations of
the obstruction class of a particular lift.
\end{statement}
\begin{proof}
We construct the free Lie bracket matrices over $\F_p$,
take the kernel from B.10, and construct the transposition
$S$ and cycle $R$ by substitution in Lyndon brackets.
The code checks $S^2=R^3=1$ and $SRS=R^{-1}$ on each
module. For a normalized cochain $c$ the differential is
\[
 (dc)(g_1,\ldots,g_{j+1})=
 g_1c(g_2,\ldots,g_{j+1})
 +\sum_{i=1}^j(-1)^i
  c(g_1,\ldots,g_i g_{i+1},\ldots,g_{j+1})
 +(-1)^{j+1}c(g_1,\ldots,g_j).
\]
Cochains with an identity argument are zero. We construct
the full matrices $d^0,d^1,d^2$, of target dimensions
$5\dim M_n$, $25\dim M_n$, and $125\dim M_n$.
Exact finite-field elimination gives
$\dim H^1=5\dim M_n-\rank d^1-\rank d^0$ and
$\dim H^2=25\dim M_n-\rank d^2-\rank d^1$.
Both consecutive matrix products are checked to be zero.
The module matrices, ranks and answers are retained in
\texttt{b\_modcohom\_matrices.json} and
\texttt{b\_modcohom.json}. This proves the displayed
finite computations independently of semisimplicity.
\end{proof}

\begin{statement}{B.14}{Open}{The exceptional-prime integral lifting problem}
The table in B.13 does not decide whether the specified
$q_7$, $2q_7$, or $3q_7$ lifts to an $S_3$-centralizing
automorphism of the free pro-$2$ or pro-$3$ group.
This return has not established those infinite lifts
or an obstruction ruling them out.
\end{statement}
\paragraph{Attacks and the exact missing calculation.}
B.11 does lift each vector without the symmetry condition.
The first attempted symmetry argument, coprime
nonabelian $H^1$ vanishing, fails because $6$ is not
invertible. We replaced this argument by the full
normalized-bar calculation B.13, which exhibits the
nonzero obstruction modules rather than suppressing
them. The supplied rational class-twelve jet was also
examined: its displayed weight-nine and weight-eleven
coefficients have denominators divisible by $2$ and
$3$ (common denominators $720$ and $8640$).
Those rational coordinates alone neither give an
integral group element nor prove that a different
integral jet cannot exist. Multiplying its first
term by a prime is not a proof that every subsequent
obstruction is killed.

We then solved the genuine integral Hall-coordinate
equations simultaneously, preserving all free parameters
from earlier degrees. The resulting integral group word
exists through class twelve, without scaling $q_7$;
the next statement gives the complete certificate.
The missing step is a compatible lift through all later
classes, or a specific nonzero obstruction beyond this
bound. A finite jet, even an integral one, does not
justify that induction.

% Proof-ready supplement for Part B. The parent may renumber B.14a.
\begin{statement}{B.14a}{Proved}{An integral group lift of the exceptional-prime witness through class twelve}
Let $F=\langle x,y\rangle$ be the discrete free group,
$N_{12}=F/\gamma_{13}F$, and $z=(xy)^{-1}$.
There is an element $f\in\gamma_7N_{12}$ whose image in
$\gamma_7F/\gamma_8F$ is the specific integral vector
$q_7$ of B.12 and which satisfies the genuine group equations
\[
 f(y,x)f(x,y)=1,\qquad
 f(z,x)f(y,z)f(x,y)=1
 \quad\hbox{in }N_{12}.
\]
The same assertion holds with $q_7$ replaced by any integral
multiple. Consequently these finite nilpotent jets exist over
$\mathbb Z_2$ and $\mathbb Z_3$ without scaling $q_7$.
This is a finite-stage assertion; it does not assert an infinite
pro-$2$, pro-$3$, or full profinite lift.
\end{statement}
\begin{proof}
We use actual group commutators $[g,h]=ghg^{-1}h^{-1}$.
For a Lyndon word $w$, let $C_x=x$, $C_y=y$, and
$C_w=[C_u,C_v]$, where $w=uv$ is the standard Lyndon
factorization with $v$ the longest proper Lyndon suffix.
The leading Magnus term of $C_w-1$ is the standard integral
Lie bracket $b_w$. These leading terms form a basis of each
$\gamma_nF/\gamma_{n+1}F$. Induction over the lower central
series therefore gives unique integral coordinates for
$\gamma_7F/\gamma_{13}F$ in the ordered products of
$C_w$, $7\leq |w|\leq12$.

Set $x=1+X$, $y=1+Y$ in the truncated noncommutative
power-series ring over $\mathbb Z$. In particular
\[
 z=(1+Y)^{-1}(1+X)^{-1},
\]
and this full group series, rather than the linear expression
$-X-Y$ or logarithmic BCH coordinates, is used in every
substitution. If $A=g-1$, $B=h-1$, the exact identity
\[
 [g,h]-1=(AB-BA)(1+A)^{-1}(1+B)^{-1}
\]
constructs every needed Hall series by finite integer arithmetic.
The product of two terms of degree at least seven has degree at
least fourteen. Hence through degree twelve the Magnus expansion
of $\prod C_w^{a_w}$ is exactly
$1+\sum a_w(C_w-1)$, and both displayed group equations
are linear equations in all the Hall exponents.

We fix the seven weight-seven exponents to the coefficients of
$q_7$ and include \emph{every} exponent in weights eight through
twelve as an unknown. The simultaneous systems, retaining all
earlier free parameters, have the following exact ranks:
\[
\begin{array}{c|rrrrr}
\text{class}&8&9&10&11&12\\\hline
\text{unknown exponents}&30&86&185&371&706\\
\text{rank}&20&66&147&302&578\\
\text{free parameters}&10&20&38&69&128.
\end{array}
\]
These ranks and the final solutions are computed by
\texttt{code/e\_exceptional\_jets.py}; its full output and exact
coordinates are in \texttt{e\_exceptional\_jets\_12.json}.
At class twelve write the integral system as $Aa=b$.
Exact elimination at the primes $2,3,7$ gives solutions
$a_i=U_i/d_i$, with $U_i$ integral and
\[
 (d_2,d_3,d_7)
 =(43641883694233125,\ 720545000,\ 24).
\]
The corresponding denominator is a unit at the prime used in
each elimination. More strongly, the explicit identity
\[
 265181d_2-16061451206962d_3-1276d_7=1
\]
gives an integral solution
\[
 a=265181U_2-16061451206962U_3-1276U_7,
 \qquad Aa=b.
\]
The last stage's \texttt{combined\_jet.hall\_coordinates}
contains this vector. Together with the seven fixed coefficients
of $q_7$, it explicitly specifies the desired group element.

The program checks every associative-word coefficient of both
group defects through degree twelve, with exact rational or
integer arithmetic. Vanishing in the truncated Magnus ring
implies the asserted identities in $N_{12}$: if a group word
had a first nonzero lower-central component of degree at most
twelve, its nonzero Lyndon Lie polynomial would be its first
nonzero Magnus term. Conversely, a word in $\gamma_{13}F$
has no such term. This also justifies the independent full-word
verification of the recorded integral solution.

For completeness, selecting Lyndon-word rows when solving does
not discard an equation. Each cycle defect is primitive, up to
the given truncation, for the Magnus Hopf structure
$\Delta X=X\otimes1+1\otimes X+X\otimes X$.
Its first nonzero homogeneous part would be an ordinary free
Lie polynomial. The triangular Lyndon basis detects that part
on Lyndon words. Induction on degree proves that vanishing of
all selected homogeneous Lyndon coefficients is equivalent to
vanishing of the entire defect. In addition the code checks the
complete associative expansions of each final solution directly.

Finally, in this degree range replacing every Hall exponent by
its multiple $m$ multiplies both linear defects by $m$ and
replaces the leading term by $m q_7$. This proves the assertion
for all integral multiples, including the two multiples singled
out in the prompt. Passing to the pro-$p$ completion of $N_{12}$
or any finite $p$-group quotient preserves the verified equations.
\end{proof}

\paragraph{What remains of the exceptional-prime problem.}
The calculation evaluates the actual group-equation lifting
problem through the stated nilpotency class; it improves on the
mere existence of nonzero cohomological obstruction spaces and
on rational BCH jets. It does not prove that one of these jets
extends through every subsequent class. Existence at all classes,
with integral $p$-adic coordinates, would suffice: compactness of
the product of the Hall-coordinate copies of $\mathbb Z_p$
would give a compatible inverse-limit solution. The required
nonempty sets have been checked only to the finite bound above.
In particular the factor $6$ in an averaging argument is not a
unit for the two primes in question; taking norms may multiply
the first term at each new stage and therefore does not establish
an infinite lift of $q_7$, $2q_7$, or $3q_7$.


\subsection{Geometric images, normalizers, and filling}

\begin{statement}{B.15}{Refuted}{A Lie algebra alone loses the tripod symmetry}
The positive geometric Lie algebra at type $(0,3)$ is zero.
Its Lie normalizer is therefore the entire special outer
algebra, not the special $S_3$-invariants. The corrected
normalizer statement retains the finite geometric $S_3$:
the Lie algebra of its normalizer in the special
automorphism group is its invariant algebra, the fourth
column of B.10. The rational graded arithmetic image is
the odd free Lie algebra in the final column.
\end{statement}
\begin{proof}
The pure mapping class group of the labelled tripod is
trivial. For any Lie algebra $\mathfrak a$,
$N_{\mathfrak a}(0)=\mathfrak a$. At weight five B.10
gives dimensions $3$ and $1$, an explicit discrepancy.
The Lie algebra of the normalizer of a finite subgroup in
characteristic zero centralizes it: infinitesimal
conjugation cannot vary an element of a finite étale
scheme. Its Lie algebra is consequently the fixed space.
B.12 realizes these fixed spaces as leading terms when
$\ell\ge5$; over $\Q_\ell$ they also describe the
prounipotent centralizer.

For the arithmetic assertion we use two distinct checked
inputs. \citet[Theorem 8.4, Theorem 7.3 and Lemma 8.5]{HainMatsumoto2003}
identify the rational graded Galois image with a quotient
of the free weighted-completion Lie algebra, and show
that it is generated by the odd Soulé classes of degrees
at least three; its degree-one generator acts trivially
on the tripod. \citet[Theorem 1.1 and the deduction of
Conjecture 1 in the introduction]{Brown2012} proves the
faithfulness of the mixed-Tate motivic action on the
tripod groupoid. The étale realization comparison used
by Hain--Matsumoto identifies its generator spaces with
the same Soulé classes. Thus a relation among the
arithmetic derivations would give a relation in that
faithful free mixed-Tate Lie algebra. No such relation
exists. Passing to outer derivations loses none of this
action: in the normalized representative of B.12 a
nonzero inner action cannot have conjugator in $F'$.
This proves equality with the odd free image, without
assuming that the full pentagon solution Lie algebra
equals it in every weight.
\end{proof}

\begin{statement}{B.16}{Proved}{The four-puncture calculations through degree six}
Write $\mathfrak p$ for the point-pushing free Lie
subalgebra and $\mathfrak n=N_{\mathfrak s_3}(\mathfrak p)$.
All dimensions in this table are exact:
\[
\begin{array}{c|rrrrrr}
n&1&2&3&4&5&6\\ \hline
\mathfrak s_{3,n}&2&1&6&6&28&36\\
\mathfrak s_{3,n}^{S_4}&0&0&1&0&3&1\\
\mathfrak p_n&2&1&2&3&6&9\\
\mathfrak n_n&2&1&3&3&7&9\\
(\mathfrak n/\mathfrak p)_n&0&0&1&0&1&0\\
\text{arithmetic image modulo geometry}&0&0&1&0&1&0\\
\text{filling image}&0&0&1&0&3&0\\
\text{filling kernel}&2&1&5&6&25&36
\end{array}
\]
In this range the normalizer quotient, the arithmetic
quotient and $\grt_1$ agree, via their arithmetic
comparison maps.
\end{statement}
\begin{proof}
In $L(x,y,z)$ choose the degree-one point pushes
\[
 u=([y,x],[x,y],0),\qquad
 v=([z,x],0,[x,z]).
\]
Their action is faithful, as follows either by the
configuration-space fibration or by the triangular
free-fibre action: successive brackets have the
independent Lyndon leading words in the two moving
paths. The code also checks their ranks independently
in every degree used. A derivation $D$ normalizes
$\mathfrak p$ if and only if both $[D,u]$ and $[D,v]$
belong to $\mathfrak p$, since Jacobi then proves the
condition on all Lie words in $u,v$. In each degree
the program constructs a basis of B.10, computes both
brackets, projects to a complementary row space of
$\mathfrak p_{n+1}$, and takes the simultaneous kernel.
This is a sufficient normalizer test, not just a
necessary one. Exact rational elimination gives the
fourth and fifth rows. The $S_4$ row is independently
obtained by the Witt character average.

Filling $x$ means substitution $x=0$ followed by passage
to outer derivations. The program computes its matrix
and its rank. An arbitrary sum-preserving tripod
derivation extends by $D(x)=0$, so filling is onto
the full tripod algebra in every positive degree.
The formula for point pushes shows that they fill
to zero. Direct kernel calculation on the computed
normalizer gives precisely $\mathfrak p$ through six.

The arithmetic action on the configuration fibration
normalizes its geometric monodromy. Restriction to the
filled tripod realizes the nonzero Soulé classes in
degrees three and five from B.15, giving the matching
lower bounds. The two cycle and special equations
bound $\grt_1$ by the invariant column of B.10, which
has the same dimensions through six. The arithmetic
injection into $\grt_1$ therefore identifies all three
graded spaces in this range. The only possible
positive-degree bracket within it is the self-bracket
of the degree-three line, which is zero. The
identification consequently respects every bracket
in the stated range.
\end{proof}

\begin{statement}{B.17}{Refuted}{Exactness onto the entire lower normalizer}
Filling $(0,4)\to(0,3)$ is not an exact sequence with
point-pushing kernel on all special automorphisms or
all special Lie algebras. On the point-pushing
normalizer the corrected continuous group sequence is
\[
 1\longrightarrow F_2^\Sigma\longrightarrow
 N_{O_{0,4,x}^\Sigma}(F_2^\Sigma)
 \longrightarrow \GT^\Sigma\times S_3\longrightarrow1,
\]
where $\Sigma$ is all primes or one prime, the
distinguished fourth cusp is fixed, and the last group
is regarded as its subgroup of outer automorphisms of
the tripod. For $\Sigma=\{\ell\}$, $\ell\ge5$, this
image is strictly smaller than the entire lower
normalizer $\GT_0^{(\ell)}\times S_3$.
\end{statement}
\begin{proof}
An explicit extra Lie-kernel element in degree three is
\[
 E(x)=[[y,[x,y]],x],\quad
 E(y)=-[[x,[x,y]],y],\quad E(z)=0.
\]
Jacobi gives $E(x+y+z)=0$. Filling $x$ gives zero;
filling $z$ gives a nonzero special outer tripod
derivation. Point-pushing fills to zero for either
choice, so $E$ is not point-pushing. Its ranks and
the two fillings are checked by the same integer
Lyndon arithmetic underlying B.16. In particular the
degree-three full kernel has dimension five whereas
point-pushing has dimension two. It is a vector-space
comparison; point-pushing is not an ideal in the full
special algebra.

We give the reduction of the normalizer sequence to
the precise configuration theorems. Put
$T=\mathbb P^1-\{0,1,\infty\}$ and
$\mathcal E=\pi_1^\Sigma(T_2)=\pi_1^\Sigma(M_{0,5})$.
The ordered two-point configuration fibration gives
$1\to F_3^\Sigma\to\mathcal E\to F_2^\Sigma\to1$.
Since $F_3^\Sigma$ is center-free, this extension is
the pullback of
$\Aut(F_3^\Sigma)\to\Out(F_3^\Sigma)$ along its
faithful monodromy subgroup. Conjugation by a
normalizing outer automorphism extends it to
$\mathcal E$, uniquely modulo inner automorphisms
from the fibre. Conversely an automorphism preserving
that fibre restricts to a normalizer element.
The extension is inner exactly for point-pushing.

\citet[Corollaries 2.6(i), 2.8(i) and Definition 2.7]{HoshiMinamideMochizuki2022}
identify $\Out(\mathcal E)=\GT^\Sigma\times S_5$
and recover the five generalized fibre subgroups
intrinsically. Fixing our fibre fixes one moving
label; preserving its distinguished diagonal cusp
fixes the other. The remaining permutation group is
$S_3$. \citet[Theorem 6.1, $n=1$]{HoshiMochizuki2011}
proves that restriction to the base is injective on
the fibre-and-cusp-preserving outer group. Thus the
kernel of filling on this normalizer is exactly
point-pushing, and the image is the displayed
$\GT^\Sigma\times S_3$. No assertion that
$\GT^\Sigma$ is itself a pro-$\Sigma$ group is needed.
Finally B.12 gives a lower-tripod element with
nonzero pentagon defect, so it is not in this image.
This is the required correction to exactness onto
the \emph{entire} lower normalizer.
\end{proof}

\begin{statement}{B.18}{Known}{The elliptic monodromy and its computed normalizer}
For the universal once-punctured elliptic curve let
$L=L(a,b)$ and $\omega=[a,b]$. Special outer
derivations have
\[
 \dim\mathfrak e_n=2\ell_{n+1}(2)-\ell_{n+2}(2)
                    -\delta_{n,2}.
\]
Its geometric graded image is
$\mathfrak{sl}_2\ltimes\mathfrak u$, where
$\mathfrak u$ is generated by the $\SL_2$-orbits of
$\epsilon_{2m}$, $m\ge2$. Write
$\mathfrak n_E^+=N_{\mathfrak e}(\mathfrak{sl}_2
\ltimes\mathfrak u)$ in positive degrees. Then:
\[
\begin{array}{c|rrrrrrrrrrrr}
n&1&2&3&4&5&6&7&8&9&10&11&12\\ \hline
\mathfrak e_n&0&0&0&3&0&6&4&13&12&37&40&99\\
\mathfrak e_n^{\langle-I\rangle}
 &0&0&0&3&0&6&0&13&0&37&0&99\\
\mathfrak u_n&0&0&0&3&0&5&0&10&0&24&0&50\\
\mathfrak n_{E,n}^+&0&0&0&3&0&6&0&10&0&25&0&50\\
(\mathfrak n_E^+/\mathfrak u)_n
 &0&0&0&0&0&1&0&0&0&1&0&0\\
\text{arithmetic quotient}
 &0&0&0&0&0&1&0&0&0&1&0&0
\end{array}
\]
The finite symmetry shown is the central involution
present on every elliptic fibre. The larger CM
symmetries depend on the fibre and are treated in G.
\end{statement}
\begin{proof}
Subtract an inner derivation to make $D(\omega)=0$.
The map
\[
 L_{n+1}^{\,2}\longrightarrow L_{n+2},\qquad
 (u,v)\longmapsto[u,b]+[a,v]
\]
is onto. The only inner derivation still preserving
$\omega$ comes from its centralizer $\Q\omega$,
so exactly one dimension is removed in degree two.
The centralizer assertion follows from the leading
word argument in the free associative algebra
(or from the centralizer of a nonpower free Lie
element). The involution acts as $(-1)^n$.

The literature reduction is
\citet[Theorem 10.7, Proposition 20.1,
Corollary 20.2, Theorem 22.3, Corollary 23.3]{HainMatsumoto2020}.
These identify the relative-completion monodromy
image and its generators. Their ordinary $W$-degree
is bracket length; their different filtration $M$
must not be substituted for it. Remark 22.2 identifies
$\epsilon_2=\ad\omega$ as the tangent-vector
generator, which disappears in the outer image.
For a reproducible normalization put $B_j=\ad_a^j b$.
Their equation (22.3), in our basis, is
\[
\begin{split}
 \epsilon_k(a)&=[B_{k-1},a],\\
 \epsilon_k(b)&=[B_{k-1},b]
 -\sum_{\substack{j+h=k-1\\j>h>0}}(-1)^j[B_j,B_h].
\end{split}
\]
The program checks $D(\omega)=0$, generates the
entire $\SL_2$-orbit by repeated lowering, forms
every bracket of lower degrees, and performs exact
rational row reduction. This recursion is exhaustive
for the Lie algebra generated by these orbits.

To prove the normalizer upper bound, semisimplicity
of finite-dimensional $\mathfrak{sl}_2$-modules gives
\[
 (\mathfrak n_E^+/\mathfrak u)_n
 \hookrightarrow(\mathfrak e_n/\mathfrak u_n)^{\mathfrak{sl}_2}.
\]
Witt characters give the ambient decomposition;
the computed homogeneous geometric basis gives the
other decomposition. The bounds are zero except
one in degree six and three in degree ten.
In degree ten impose the further necessary condition
$[D,\epsilon_4]\in\mathfrak u_{14}$. The program
computes $\dim\mathfrak u_{14}=109$; adjoining the
three invariant brackets raises its rank to $111$.
The resulting degree-ten upper bound is one.
For the matching lower bounds,
\citet[Proposition 29.1 and Remark 29.3]{HainMatsumoto2020}
construct an injective arithmetic mixed-Tate map to
this normalizer quotient and place $\sigma_{2m+1}$
in ordinary degree $4m+2$ modulo geometry.
It supplies the degree-six and degree-ten lines.
The only inner line removed was in degree two,
so neither arithmetic class disappears. This proves
every entry and the asserted arithmetic equality
through twelve.
\end{proof}

\begin{statement}{B.19}{Refuted}{The epsilon algebra is not the normalizer quotient}
At type $(1,1)$ the proposed identification of the
normalizer \emph{modulo geometry} with the algebra
generated by the $\epsilon_{2k}$ is false. In the
range of B.18 the corrected quotient is the odd
mixed-Tate Lie algebra with degrees doubled.
\end{statement}
\begin{proof}
The nonzero three-dimensional $\SL_2$-orbit of
$\epsilon_4$ is geometric and is zero in the
quotient. The quotient itself is zero in degree
four, and has its first nonzero class in degree six.
This is already a dimension counterexample.
Through ordinary degree twelve its two lines
are $\sigma_3,\sigma_5$ in degrees six and ten.
Their brackets cannot produce a new class in the
stated range: the self-bracket in degree twelve is
zero. Pollack's relations constrain the geometric
epsilon algebra. They cannot turn that geometric
algebra into a quotient in which it has been killed.
An all-degree equality of the elliptic normalizer
quotient with $\grt_1$ is not a consequence of the
finite table or of the comparison theorem.
\end{proof}

\subsection{Universal properties and comparison maps}

\begin{statement}{B.20}{Proved}{The integral envelope of each virtually abelian stage}
Let $\Gamma$ be a finitely generated group with a normal lattice
$M\cong\Z^r$ of finite index. For $n\ge1$ form the group scheme
$\Gamma_{nM}$ whose components are indexed by $\Gamma/nM$ and whose
identity component is $(nM)\otimes\Ga$. Put
\[
 E_\Gamma=\varprojlim_{n\mid n'}\Gamma_{nM}.
\]
For every ring $A$ and every closed finitely presented subgroup
$G\subset\GL_{d,A}$, restriction to the canonical homomorphism
$\Gamma\to E_\Gamma(A)$ induces a bijection
\[
 \Hom_{A\text{-groups}}((E_\Gamma)_A,G)
 \ \cong\
 \{\rho:\Gamma\to G(A):\rho(\gamma)
       \text{ is fibrewise quasi-unipotent for all }\gamma\in\Gamma\}.
\]
The construction is independent of $M$ up to the unique isomorphism
compatible with $\Gamma$. This is the functorial, direct-limit reading
of \MS{50}{663--663'}{31--32}, rather than an inverse limit of Hom sets.
\end{statement}
\begin{proof}
Choose coset representatives $t_h$ for $\Gamma/nM$. In lattice coordinates
the product has the integral polynomial formula
\[
 (v,h)(w,k)=(v+h\cdot w+c(h,k),hk),
 \qquad c(h,k)=t_ht_kt_{hk}^{-1}\in nM.
\]
Associativity is exactly the cocycle identity satisfied by these actual
representatives. Replacing representatives translates coordinates by
integral vectors. Refinement of lattices gives homomorphisms of these
schemes, hence the stated pro-group.

Since $G$ is affine and finitely presented, a map from the inverse limit
factors through some $\Gamma_{nM,A}$; its group identities also hold at
a sufficiently late stage. Thus the Hom set on the left is
$\varinjlim_n\Hom(\Gamma_{nM,A},G)$. Every element of $\Gamma$ has
a positive power in $nM$. The image of $nM$ under a stage morphism
is contained in a commuting product of additive one-parameter groups.
In every characteristic-zero fibre those images are unipotent;
in characteristic $p$ an additive-group image has exponent $p$.
This proves fibrewise quasi-unipotence.

Conversely let $\rho$ have the stated property and choose a basis of $M$.
The proof of D.2 supplies a common positive multiple of that basis whose
images extend to additive one-parameter groups in $G$. Their matrix
formulas are polynomials in the commuting matrices $\rho(m_i)-1$
after a preliminary common power; consequently the one-parameter
groups commute. They define an additive-group homomorphism on the
chosen lattice coordinates. Conjugation by each of the finitely many
coset representatives, and each of their multiplication relations,
give a finite set of polynomial identities. These identities hold
on all integer lattice points.

The torsion correction is essential here. If a polynomial over $A$
vanishes on every point of $\Z^r$, repeated univariate Vandermonde
evaluation gives a nonzero integer annihilating all its coefficients.
For identities comparing two maps, the constant coefficient already
vanishes. Replacing every variable by $Dv_i$, where $D$ annihilates
the other coefficients, therefore makes the identities hold over $A$.
Only finitely many coefficient equations occur, because the lattice,
component group, and equations of $G$ are finite. A common dilation
gives a stage $\Gamma_{nM,A}\to G$. On a component use its actual
representative in $\Gamma$ and the additive map on $nM$; the checked
conjugation and cocycle identities prove it is a group morphism.

The same argument applied to two stage morphisms agreeing on $\Gamma$
makes them agree after a common lattice refinement. This proves
injectivity in the direct limit. Two normal lattices have a common
finite-index sublattice, proving independence of the choice. This
reproduces the proof behind LM1 E.14 rather than adopting its label
as authority.
\end{proof}

\begin{statement}{B.21}{Proved}{Filling and the virtually abelian system}
Let $F=F_r$ and let $q:F\twoheadrightarrow F_{r-1}$ kill one chosen
free generator, with kernel $R$. The induced continuous map
$V(q):V(F)\to V(F_{r-1})$ is split surjective and
\[
 \ker V(q)=
 \varprojlim_{K\triangleleft_f F} R/(R\cap[K,K]).
\]
The same split surjectivity holds for the inverse systems of
\S50 envelopes and for every nilpotent Malcev truncation.
For the latter the kernel Lie algebra is the ideal generated by
the killed generator, modulo brackets exceeding the chosen class.
These statements do not identify the kernel with the absolute
completion of the infinitely generated $R$ without a separate
cofinality assertion.
\end{statement}
\begin{proof}
The homomorphism inserting the remaining generators is a section of
$q$, so functoriality supplies sections for every construction.
The finite-index normal subgroups $q(K)$ are cofinal in $F_{r-1}$:
given $J\triangleleft_fF_{r-1}$, take $K=q^{-1}J$.
For each $K$,
\[
 \ker\{F/[K,K]\to F_{r-1}/[q(K),q(K)]\}
 =R[K,K]/[K,K].
\]
Indeed $q([K,K])=[q(K),q(K)]$, so an element in the kernel can
be corrected by an element of $[K,K]$ to lie in $R$.
Taking limits proves the kernel formula. The componentwise
polynomial maps of B.20 prove the envelope assertion; the
section remains a section. Free Lie algebras have the analogous
universal quotient property, proving the Malcev assertion.
Here $R$ has infinite index in $F$; the cofinality proof for a
finite cover in B.5 does not apply.
\end{proof}

\begin{statement}{B.22}{Proved}{Precise comparisons with weighted and Teichm\"uller completions}
There is a natural morphism from the pro-object $\mathscr E$ of B.2
to the image closure of every compatible adelic finite-jet
representation furnished by $\GT$ or $\GT_{\rm ell}$. There is no
unconditional identification of these targets with $\mathscr E$.

More precisely, fix a continuous Zariski-dense reductive representation
$\rho_\ell:\mathcal M\to R(\Q_\ell)$ and a central cocharacter
$w:\Gm\to R$. Its weighted completion
$\mathcal W_{\rho_\ell,w}$ maps uniquely to a finite-type target
$H_\ell$ over $R$ whenever the representation into $H_\ell(\Q_\ell)$
lifts $\rho_\ell$ and $\ker(H_\ell\to R)$ is unipotent with
strictly negative $w$-weights on its abelianization. For targets
defined over $\Q$ with the compatible adelic representation of B.1,
the comparison diagram is
\[
 \mathscr E_{\Q_\ell}\longrightarrow (H_{\mathbb A})_{\Q_\ell}
 \longleftarrow\mathcal W_{\rho_\ell,w},
\]
where $H_{\mathbb A}$ is the $\Q$-Zariski closure of the adelic
image. The right arrow has image the $\Q_\ell$-Zariski closure of
the $\ell$-adic image; these two closures need not be equal.
\end{statement}
\begin{proof}
For a prescribed system of finite-jet targets take the $\Q$-Zariski
closure of the adelic image in each target. If its character and
rationality on $\mathcal G_0$ are as in B.1, each closure is an
object of that category, so projection from its inverse limit gives
the left arrow. Compatibility under truncation gives a morphism
of pro-groups. This supplies a construction, not an assertion that
an arbitrary completion has the required rational diagonal action.

The second assertion is the universal property in
\citet[Definition 4.1 and Proposition 4.4]{HainMatsumoto2003}.
It also follows from the product-and-image-closure construction
of B.2 restricted to negatively weighted unipotent extensions
of the fixed $R$. The canonical $\mathcal M$-image is Zariski
dense in weighted completion, so the scheme-theoretic image is
exactly the closure of the given representation. Negative
weights and a fixed reductive representation are extra
hypotheses, not imposed by B.1.

For the genus-zero arithmetic tower use its compatible
representation of $\Gamma_\Q$ into $\widehat\GT$, followed
prime by prime by unipotent finite jets. It is trivial on
$\mathcal G_0$ in the arithmetic quotient. Its character is
cyclotomic, so the preceding construction applies to that
quotient. Choosing an associator identifies $\GT$ with its
graded form $\GRT$; the identification depends on that choice.
The negative-weight property constructs a map from cyclotomic
weighted completion to the arithmetic image, not a surjection
onto all of $\GT$ by definition.

For the elliptic tower use the morphism from the arithmetic
fundamental group of the tangent-point moduli stack to
$\widehat\GT_{\rm ell}$ of
\citet[Theorem 8.4 and Proposition 8.5]{Enriquez2014}.
The geometric $B_3$ action has rational word formulas, so
each adelic algebraic jet satisfies the same rationality check.
The forgetful morphism $\GT_{\rm ell}\to\GT$ is split by
\citet[Proposition 2.20]{Enriquez2014}. With his convention
$c=(Y,X)=YXY^{-1}X^{-1}$ the section is
\[
\begin{split}
 g_+(X,Y)&=f(X,c)X^\lambda f(X,c)^{-1},\\
 g_-(X,Y)&=c^{(\lambda-1)/2}
      f(YX^{-1}Y^{-1},c)\,Y\,f(X,c)^{-1}.
\end{split}
\]
His proposition verifies the elliptic braid identities and
multiplicativity and gives the same formulas for the
proalgebraic group-scheme version. Projection recovers
$(\lambda,f)$, proving it is a section. Its kernel is therefore
additional elliptic data, containing the geometric action;
it is not the arithmetic quotient in B.18. These are the
typed comparisons requested in Lead 1. A single equality
cannot follow from the phrase \emph{Lie approximation}.
\end{proof}

\subsection{The exhaustive finite-cover computation}

\begin{statement}{B.23}{Proved}{Every finite quotient of order at most twenty-four}
The program enumerates all $74$ isomorphism types of groups
of order at most $24$ and all $9714$ epimorphisms from a
fixed based $F_2$ to chosen representatives. Exactly $65$
isomorphism types are two-generated. For each epimorphism
it computes a Schreier basis of $N$, the complete matrix
$N^{\ab}\to F_2^{\ab}$, its image lattice, a presentation of
the derived module of $N/N''$, and the induced graded maps
through lower-central class four. Every exact word and
lattice check succeeds.
\end{statement}
\begin{proof}
We use GAP 4.15.1 and SmallGrp 1.7.0, the complete
irredundant catalogue in this range \citep{SmallGrp2025}.
For each group the program tests every ordered pair and
retains it exactly when it generates the whole group.
Such a pair is exactly an epimorphism from the specified
based free group, so none are lost by quotienting by
Nielsen moves or automorphisms of the finite target.

Breadth-first traversal of the positive Cayley graph gives a
spanning tree. Each non-tree directed edge yields a reduced
Schreier word. There are $2d-(d-1)=d+1$ words; evaluation in
the multiplication table checks that each is in $N$.
The graph's cellular chain complex shows these words are a
free basis. Their exponent vectors give the $2$-row matrix.
Its column lattice has index $|G^{\ab}|$, as follows
independently from $F_2^{\ab}/\im N^{\ab}=G^{\ab}$;
the program checks this equality for every pair.

For $r=d+1$, put $R_r=\Z[X_1^{\pm1},\ldots,X_r^{\pm1}]$.
The universal abelian cover of a rose with $r$ loops has no
two-cells, so its cellular boundary gives the exact module
description
\[
 N'/N''=\ker(R_r^r\longrightarrow R_r),
 \qquad e_i\longmapsto X_i-1.
\]
Thus the metabelian group's abelianization has rank $r$
and its derived module has rank $r-1$ over the fraction field.
The presentation is recorded for each actual Schreier basis;
the module is not replaced by its rank.

Lower-central quotients of a free group have the integral
Lyndon basis. If the degree-one image lattice has Smith
invariants $s_1\mid s_2$, the image Smith invariants in
degrees $1,2,3,4$ are
\[
\begin{gathered}
 (s_1,s_2),\qquad(s_1s_2),\qquad
 (s_1^2s_2,s_1s_2^2),\\
 (s_1^3s_2,s_1^2s_2^2,s_1s_2^3).
\end{gathered}
\]
Each indicated multidegree is a free rank-one summand;
an integral change of bases induces an integral Lie
automorphism. Kernel ranks are $\ell_j(r)-\ell_j(2)$.
These are the graded maps and free nilpotent presentations,
not an inference of filtered injectivity from a graded rank.
The comparison kernels and images themselves are B.7.
Comparisons are flagged undefined unless $G$ is in the
required variety of groups.

The table prints every isomorphism type, with SmallGrp
index, number $e$ of epimorphisms, $|G^{\ab}|$, derived
length $d_{\rm sol}$, and nilpotence class $c$ (a dash means
nonnilpotent). For $e>0$ the Schreier rank is $|G|+1$;
individual matrices and words are in the machine-readable
output. Cases with $e=0$ are retained explicitly.
\end{proof}
{\small
\begin{longtable}{rrlrrrr}
\toprule $|G|$&ID&Structure&$e$&$|G^{\rm ab}|$&$d_{\rm sol}$&$c$\\\midrule
\endfirsthead
\toprule $|G|$&ID&Structure&$e$&$|G^{\rm ab}|$&$d_{\rm sol}$&$c$\\\midrule
\endhead
\bottomrule\endfoot
1 & 1 & 1 & 1 & 1 & 0 & 0\\
2 & 1 & C2 & 3 & 2 & 1 & 1\\
3 & 1 & C3 & 8 & 3 & 1 & 1\\
4 & 1 & C4 & 12 & 4 & 1 & 1\\
4 & 2 & C2 $\times$ C2 & 6 & 4 & 1 & 1\\
5 & 1 & C5 & 24 & 5 & 1 & 1\\
6 & 1 & S3 & 18 & 2 & 2 & --\\
6 & 2 & C6 & 24 & 6 & 1 & 1\\
7 & 1 & C7 & 48 & 7 & 1 & 1\\
8 & 1 & C8 & 48 & 8 & 1 & 1\\
8 & 2 & C4 $\times$ C2 & 24 & 8 & 1 & 1\\
8 & 3 & D8 & 24 & 4 & 2 & 2\\
8 & 4 & Q8 & 24 & 4 & 2 & 2\\
8 & 5 & C2 $\times$ C2 $\times$ C2 & 0 & 8 & 1 & 1\\
9 & 1 & C9 & 72 & 9 & 1 & 1\\
9 & 2 & C3 $\times$ C3 & 48 & 9 & 1 & 1\\
10 & 1 & D10 & 60 & 2 & 2 & --\\
10 & 2 & C10 & 72 & 10 & 1 & 1\\
11 & 1 & C11 & 120 & 11 & 1 & 1\\
12 & 1 & C3 $\rtimes$ C4 & 72 & 4 & 2 & --\\
12 & 2 & C12 & 96 & 12 & 1 & 1\\
12 & 3 & A4 & 96 & 3 & 2 & --\\
12 & 4 & D12 & 36 & 4 & 2 & --\\
12 & 5 & C6 $\times$ C2 & 48 & 12 & 1 & 1\\
13 & 1 & C13 & 168 & 13 & 1 & 1\\
14 & 1 & D14 & 126 & 2 & 2 & --\\
14 & 2 & C14 & 144 & 14 & 1 & 1\\
15 & 1 & C15 & 192 & 15 & 1 & 1\\
16 & 1 & C16 & 192 & 16 & 1 & 1\\
16 & 2 & C4 $\times$ C4 & 96 & 16 & 1 & 1\\
16 & 3 & (C4 $\times$ C2) $\rtimes$ C2 & 96 & 8 & 2 & 2\\
16 & 4 & C4 $\rtimes$ C4 & 96 & 8 & 2 & 2\\
16 & 5 & C8 $\times$ C2 & 96 & 16 & 1 & 1\\
16 & 6 & C8 $\rtimes$ C2 & 96 & 8 & 2 & 2\\
16 & 7 & D16 & 96 & 4 & 2 & 3\\
16 & 8 & QD16 & 96 & 4 & 2 & 3\\
16 & 9 & Q16 & 96 & 4 & 2 & 3\\
16 & 10 & C4 $\times$ C2 $\times$ C2 & 0 & 16 & 1 & 1\\
16 & 11 & C2 $\times$ D8 & 0 & 8 & 2 & 2\\
16 & 12 & C2 $\times$ Q8 & 0 & 8 & 2 & 2\\
16 & 13 & (C4 $\times$ C2) $\rtimes$ C2 & 0 & 8 & 2 & 2\\
16 & 14 & C2 $\times$ C2 $\times$ C2 $\times$ C2 & 0 & 16 & 1 & 1\\
17 & 1 & C17 & 288 & 17 & 1 & 1\\
18 & 1 & D18 & 162 & 2 & 2 & --\\
18 & 2 & C18 & 216 & 18 & 1 & 1\\
18 & 3 & C3 $\times$ S3 & 144 & 6 & 2 & --\\
18 & 4 & (C3 $\times$ C3) $\rtimes$ C2 & 0 & 2 & 2 & --\\
18 & 5 & C6 $\times$ C3 & 144 & 18 & 1 & 1\\
19 & 1 & C19 & 360 & 19 & 1 & 1\\
20 & 1 & C5 $\rtimes$ C4 & 240 & 4 & 2 & --\\
20 & 2 & C20 & 288 & 20 & 1 & 1\\
20 & 3 & C5 $\rtimes$ C4 & 240 & 4 & 2 & --\\
20 & 4 & D20 & 120 & 4 & 2 & --\\
20 & 5 & C10 $\times$ C2 & 144 & 20 & 1 & 1\\
21 & 1 & C7 $\rtimes$ C3 & 336 & 3 & 2 & --\\
21 & 2 & C21 & 384 & 21 & 1 & 1\\
22 & 1 & D22 & 330 & 2 & 2 & --\\
22 & 2 & C22 & 360 & 22 & 1 & 1\\
23 & 1 & C23 & 528 & 23 & 1 & 1\\
24 & 1 & C3 $\rtimes$ C8 & 288 & 8 & 2 & --\\
24 & 2 & C24 & 384 & 24 & 1 & 1\\
24 & 3 & SL(2,3) & 384 & 3 & 3 & --\\
24 & 4 & C3 $\rtimes$ Q8 & 144 & 4 & 2 & --\\
24 & 5 & C4 $\times$ S3 & 144 & 8 & 2 & --\\
24 & 6 & D24 & 144 & 4 & 2 & --\\
24 & 7 & C2 $\times$ (C3 $\rtimes$ C4) & 144 & 8 & 2 & --\\
24 & 8 & (C6 $\times$ C2) $\rtimes$ C2 & 144 & 4 & 2 & --\\
24 & 9 & C12 $\times$ C2 & 192 & 24 & 1 & 1\\
24 & 10 & C3 $\times$ D8 & 192 & 12 & 2 & 2\\
24 & 11 & C3 $\times$ Q8 & 192 & 12 & 2 & 2\\
24 & 12 & S4 & 216 & 2 & 3 & --\\
24 & 13 & C2 $\times$ A4 & 288 & 6 & 2 & --\\
24 & 14 & C2 $\times$ C2 $\times$ S3 & 0 & 8 & 2 & --\\
24 & 15 & C6 $\times$ C2 $\times$ C2 & 0 & 24 & 1 & 1\\
\end{longtable}}


\begin{statement}{B.24}{Open}{The characteristic-two conjugator-scheme comparison}
For the Fermat quotient $F_2\to(\Z/2)^2$, the relative completion
has radical $U_5$, whose abelianization is
$\Q\oplus\Q[(\Z/2)^2]$. The corresponding integral
characteristic-two action of the manuscript's signed
$M_{\rho,\sigma}$ and $N_{\rho,\sigma}$, with every
infinitesimal parameter, has not been constructed here.
In particular its kernel and image as group-scheme
morphisms have not been computed.
\end{statement}
The unresolved map is from the signed automorphism
functor defined by the conjugators and extra generator data
of \S49 to the automorphism functor of the \emph{relative}
Fermat completion, preserving its finite quotient and
peripheral structure. B.8 proves that replacing its radical
by the absolute $U_2$ loses three abelian dimensions already
at $n=2$. LM1 E.25, whose explicit dual-number examples
were reopened, also shows that images of two discrete
generators alone do not determine this morphism on
$\F_2[\epsilon]/(\epsilon^2)$: the extra infinitesimal
parameter must be carried along. The construction
by word substitution therefore gives a discrete-point
action only and does not determine the requested functor.
The first remaining test is that dual-number functor,
followed by compatibility with the finite quotient's
nonsemisimple action in characteristic two. The complete
finite-cover module and Schreier words are supplied by
B.23; what remains missing is the scheme action.




\subsection{The relative $S_3$ diagram for \S47}

The additional condition of \S47 concerns two different homomorphisms.
Its anchors are cote 140--3, folios 415--420, scans 229--240, especially
equations (14), (15), (20), and (30). We use the explicitly checked
dictionary in the supplied LM2 return, B.m3 and E.4, and its warning
B.g12 about finite quotients. The original cote-140--3 images were not
available for a fresh image check in this return. The definitions and
proofs below are independent algebraic statements about the displayed
maps; they do not depend on a doubtful manuscript letter.

Put
\[
 B=\langle s,r\mid s^2=r^3=1\rangle,
 \qquad p:B\longrightarrow S_3,
 \qquad K=\ker p,
\]
where $p(s)$ is a transposition and $p(r)$ a three-cycle. The group $K$
is free on $x=(rs)^2$ and $y=(sr)^2$; set $z=(xy)^{-1}$. Distinguish
the abstract free group $F_0=\langle X,Y\rangle$, $Z=(XY)^{-1}$, and
the two maps
\[
 \begin{aligned}
 \iota:F_0&\stackrel{\sim}{\longrightarrow}K\subset B,
      &(X,Y)&\longmapsto(x,y),\\
 \varphi:F_0&\twoheadrightarrow B,
      &(X,Y,Z)&\longmapsto(s,sr,r^{-1}).
 \end{aligned}
\]
Thus $v=p\varphi$ is a second quotient onto $S_3$, and
$N=\ker v$ is a different free group. Write $c:N\hookrightarrow F_0$
for inclusion and $d=\varphi|_N:N\twoheadrightarrow K$ for restriction.
All Lie algebras in this subsection are over $\Q$; degree $n$ means
that a derivation raises bracket degree by $n$.

\begin{statement}{B.25}{Proved}{The correct relative source has rank seven}
Relative completion with respect to the indicated finite quotients has
exact sequences
\[
 \begin{aligned}
 1\longrightarrow U(K)\longrightarrow\mathscr R_p(B)
       &\longrightarrow S_3\longrightarrow1,\\
 1\longrightarrow U(N)\longrightarrow\mathscr R_v(F_0)
       &\longrightarrow S_3\longrightarrow1,
 \end{aligned}
\]
where $U(K)$ and $U(N)$ are free prounipotent groups of ranks two and
seven. The quotient $\varphi$ induces
$\Phi:\mathscr R_v(F_0)\to\mathscr R_p(B)$, the identity on $S_3$,
whose map on radicals is $U(d)$. The absolute completion map induces
$\Theta:\mathscr R_v(F_0)\to U(F_0)$, whose map on radicals is $U(c)$.
The isomorphism $\iota_U:U(F_0)\to U(K)$ is a third map.

The associated graded source is the Lie algebra of an eight-punctured
sphere:
\[
 L_7=L(A_0,A_1,A_2,B_0,B_1,B_2,C_0,C_1)
       /(A_0+A_1+A_2+B_0+B_1+B_2+C_0+C_1).
\]
Writing both rank-two targets as $L_2=L(X,Y,Z)/(X+Y+Z)$ through
$\iota_U$, the maps are
\[
 \begin{array}{c|ccc}
       & A_i&B_i&C_i\\ \hline
 c     &2X&2Y&3Z\\
 d     &0&(X,Y,Z)_i&0.
 \end{array}
\]
The deck group permutes each three-element cluster $A$ and $B$ by
the standard $S_3$ action and permutes the two $C$ punctures by sign.
It acts trivially on the base target of $c$ and by puncture permutation
on the target of $d$. In degree one the joint map $(c,d)$ has rank four
and kernel isomorphic to $\mathrm{sgn}\oplus\mathrm{std}$.
\end{statement}
\begin{proof}
Relative completion means the universal prounipotent extension of the
chosen reductive quotient admitting the given homomorphism of the
discrete group; a finite group in characteristic zero is allowed
\citep[\S1, pp.~47--48]{Hain1998Relative}. For a finite quotient with
kernel $H$, its radical is $U(H)$. To see the relevant exactness,
restrict any representation of the relative completion to $H$.
Conversely, induce a unipotent representation of $H$ over the finitely
many cosets. Its algebraic image has the prescribed finite permutation
quotient and a subgroup of a product of unipotent groups as kernel.
This shows that every unipotent quotient of $H$ is detected by the
relative completion. Restriction and induction prove the universal
identification, as in B.8. Naturality gives $\Phi$ and $\Theta$.
Schreier's formula gives $\operatorname{rank}N=1+6(2-1)=7$.
The orbifold Euler characteristic of $C_2*C_3$ is $-1/6$, so its
torsion-free index-six kernel $K$ has Euler characteristic $-1$ and
rank two; Reidemeister rewriting gives the displayed basis $x,y$.

The maps can also be seen simultaneously in the rational covering
\[
 T(\lambda)=
 \frac{27\lambda^2(1-\lambda)^2}{4(1-\lambda+\lambda^2)^3}.
\]
Remove $0,1,\infty,-1,1/2,2,\zeta,\bar\zeta$ from the source,
where $\zeta=(1+\sqrt{-3})/2$, and remove $0,1,\infty$ from the
$T$-line. This is the regular degree-six cover with monodromy $v$.
The deck transformations $\lambda\mapsto1-\lambda$ and
$\lambda\mapsto1/(1-\lambda)$ generate $S_3$.
The punctures over $T=1,0,\infty$ have local degrees
$(2,2,2)$, $(2,2,2)$, and $(3,3)$, respectively. In particular,
\[
 4(1-\lambda+\lambda^2)^3-27\lambda^2(1-\lambda)^2
  =(\lambda+1)^2(2\lambda-1)^2(\lambda-2)^2.
\]
Euler characteristic gives genus zero and eight punctures.
Label the punctures over $1$ by
$(A_0,A_1,A_2)=(-1,2,1/2)$, and those over $0$ by
$(B_0,B_1,B_2)=(1,0,\infty)$. The two roots of
$1-\lambda+\lambda^2$ form the $C$ cluster. Filling the five
$A$ and $C$ punctures gives the target tripod in the $\lambda$-line.
It realizes $d$; the covering realizes $c$. Peripheral classes in
homology therefore have the displayed images. Since all three graded
Lie algebras are free on their first graded pieces, these first-degree
formulas specify the complete graded Lie homomorphisms.

The permutation representation on eight punctures has character
$(8,2,2)$ on the identity, a transposition, and a three-cycle.
Removing its diagonal line gives
\[
 H_1(N,\Q)\cong2\mathbf1\oplus\mathrm{sgn}\oplus2\mathrm{std},
 \qquad\chi=(7,1,1).
\]
In the free basis $A_0,A_1,A_2,B_0,B_1,B_2,C_0$, the two matrices are
\[
 c=\begin{pmatrix}2&2&2&0&0&0&-3\\0&0&0&2&2&2&-3\end{pmatrix},
 \qquad
 d=\begin{pmatrix}0&0&0&1&0&-1&0\\0&0&0&0&1&-1&0\end{pmatrix}.
\]
They have joint rank four, proving the last assertion.
\end{proof}

Here is a correctly typed formulation of the extra datum. For specified
automorphisms $a_v$ of $\mathscr R_v(F_0)$, $a$ of $U(F_0)$, and $b$
of $\mathscr R_p(B)$, with $a_v,b$ inducing the identity on the finite
quotients, the strict diagram asks for
\[
 \Theta a_v=a\Theta,\qquad \Phi a_v=b\Phi,\qquad
 b|_{U(K)}=\iota_Ua\iota_U^{-1}.
\]
An outer diagram allows the appropriate target inner automorphisms.
These are not equations involving $a$ alone. The absolute automorphism
$a$ does not specify $a_v$ on a rank-seven radical. In particular,
writing $U_2\to\mathscr R_p(B)$ in place of $\Phi$ would discard its
nontrivial component map and change the problem.

\begin{statement}{B.26}{Proved}{Every invariant graded action has a compatible relative lift}
Let $\mathfrak s_{r,n}$ denote degree-$n$ derivations tangential to
every peripheral class of the $(r+1)$-punctured sphere, modulo inner
derivations, as in B.10. For each $n\ge1$ the maps in B.25 induce a
surjection
\[
 R_n:\mathfrak s_{7,n}^{S_3}
       \longrightarrow
       \mathfrak s_{2,n}\oplus\mathfrak s_{2,n}^{S_3},
 \qquad [D]\longmapsto(c_*[D],d_*[D]).
\]
Consequently, for every $q\in\mathfrak s_{2,n}^{S_3}$ there is a
peripheral, deck-invariant relative graded lift whose two induced
outer actions are $(q,q)$. If $q\ne0$, there is also a lift with the
same absolute projection and induced pair $(q,0)$, which fails the
compatibility condition. Thus existential graded compatibility
imposes no additional equation on the special invariant $q$.
\end{statement}
\begin{proof}
First the map is defined. An invariant outer source class has an
invariant representative by averaging over the finite deck group.
A tangential derivation preserves the ideal generated by the $A$ and
$C$ punctures, which is the kernel of $d$. The map $c$ is the maximal
graded quotient on which the deck group acts trivially: its kernel
is the ideal generated by $gv-v$ in degree one. Indeed the
coinvariants of $H_1(N,\Q)$ are two-dimensional and $c$ is their
isomorphism onto $\Q X\oplus\Q Y$. An equivariant derivation
preserves that ideal. Its two descendants are tangential, and the
second is invariant. Changing the source representative by an inner
derivation changes both descendants by inner derivations, so $R_n$
is well-defined on outer classes.

We construct independent lifts for both factors. Write a representative
of $a\in\mathfrak s_{2,n}$ as
\[
 a(e)=[e,a_e(X,Y)],\quad e=X,Y,Z,\qquad
 \sum_{e=X,Y,Z}[e,a_e]=0.
\]
Set
\[
 X'=\tfrac16\sum_iA_i,\qquad
 Y'=\tfrac16\sum_iB_i,\qquad
 Z'=\tfrac16\sum_iC_i=-X'-Y'.
\]
Define its cabling $C(a)$ by
\[
 \begin{aligned}
 C(a)(A_i)&=[A_i,a_X(X',Y')],\\
 C(a)(B_i)&=[B_i,a_Y(X',Y')],\\
 C(a)(C_i)&=[C_i,a_Z(X',Y')].
 \end{aligned}
\]
The image of the peripheral relation is
$6\sum_e[e',a_e(X',Y')]=0$. Thus this is a tangential derivation of
$L_7$. It is deck-invariant, since every cluster sum is fixed.
Moreover $c(X')=X$, $c(Y')=Y$, and $d(X')=d(Y')=0$. Positive-degree
conjugators vanish at zero, so
\[
                 R_n(C(a))=(a,0).
\]
A common inner change of the three $a_e$ produces a source inner
change, so the construction also respects outer classes.

For the second factor take $b\in\mathfrak s_{2,n}^{S_3}$. The
rank-two outer space is zero in degree one, so assume $n\ge2$.
Average a representative of $b$ over $S_3$. In these degrees the
centralizer of a peripheral generator in the free Lie algebra has
no degree-$n$ part; hence its three conjugators $b_X,b_Y,b_Z$ are
unique and equivariant. Subtract their common mean, an inner change,
to arrange
\[
 b_X+b_Y+b_Z=0,\qquad g(b_e)=b_{ge}.
\]
Subtracting the common conjugator $b_Z$ gives a representative $b^0$
with conjugators $(b_X-b_Z,b_Y-b_Z,0)$. It fixes $Z$ and therefore
fixes $X+Y$. Substitute $X=B_0$, $Y=B_1$ in its first two images,
and set the images of $A_0,A_1,A_2,B_2,C_0,C_1$ to zero. Call the
result $E$. The identity $b^0(X+Y)=0$ proves that $E$ preserves the
eight-puncture relation; it is tangential to all eight classes.
Its $d$-descendant is $b^0$. Its $c$-descendant is zero, because
$c(B_0)=c(B_1)=2Y$ and a homogeneous Lie polynomial of degree at
least two in one letter is zero.

Finally put $T(b)=\frac16\sum_g gEg^{-1}$. It is invariant and
still has zero $c$-descendant. Its $d$-descendant is the average of
$b^0$. The averaged inner discrepancy is the inner derivation with
conjugator
$\frac16\sum_g g(b_Z)=\frac13(b_X+b_Y+b_Z)=0$.
Consequently
\[
                 R_n(T(b))=(0,b).
\]
Combining the two constructions proves surjectivity in all positive
degrees. The classes $C(q)+T(q)$ and $C(q)$ have the two pairs asserted
in the statement. When $q$ is nonzero as an outer class, their
different target projections remain different modulo inner derivations.
\end{proof}

\begin{statement}{B.27}{Proved}{Exact relative dimensions through weight nine}
Put $s_n=\dim\mathfrak s_{7,n}^{S_3}$,
$a_n=\dim\mathfrak s_{2,n}$, and
$t_n=\dim\mathfrak s_{2,n}^{S_3}$. Write $k_n=\dim\ker R_n$ and
$h_n=\dim R_n^{-1}(\operatorname{diag}\mathfrak s_{2,n}^{S_3})$.
The following dimensions are exact:
\[
\begin{array}{c|rrrrrrrrr}
 n&1&2&3&4&5&6&7&8&9\\ \hline
 s_n&4&5&38&120&693&3163&17356&92930&524166\\
 a_n&0&0&1&0&3&0&6&4&13\\
 t_n&0&0&1&0&1&0&2&1&4\\
 k_n&4&5&36&120&689&3163&17348&92925&524149\\
 h_n&4&5&37&120&690&3163&17350&92926&524153
\end{array}
\]
The image of the compatible relative space in the absolute special
invariants has dimension $t_n$. In particular the explicit $q_7$ of
B.12 passes this existential relative graded test.
\end{statement}
\begin{proof}
Let $V=H_1(N,\Q)$ and let $P$ be the permutation module on its eight
punctures. Their characters on $(1,(12),(123))$ are $(7,1,1)$ and
$(8,2,2)$. The exact special-derivation sequence in B.10 gives
\[
 \operatorname{tr}(g\mid\mathfrak s_{7,n})
 =\chi_V(g)\operatorname{tr}(g\mid L_nV)
   -\operatorname{tr}(g\mid L_{n+1}V)-\delta_{n,1}\chi_P(g),
\]
where
\[
 \operatorname{tr}(g\mid L_nV)
   =\frac1n\sum_{d\mid n}\mu(d)\chi_V(g^d)^{n/d}.
\]
Averaging over the three conjugacy classes, with sizes $1,3,2$,
gives $s_n$. The analogous rank-two formula gives $a_n$ and $t_n$.
For example the source characters in weights seven, eight, and nine
are $(103236,300,0)$, $(558404,-300,38)$, and
$(3140032,1680,-38)$. These give $17356$, $92930$, and $524166$.
Surjectivity in B.26 gives
\[
 \dim\ker R_n=s_n-a_n-t_n,\qquad
 \dim R_n^{-1}(\operatorname{diag}\mathfrak s_{2,n}^{S_3})=s_n-a_n.
\]

The executable \texttt{b\_relative47.py} independently recomputes
every rank-two special invariant basis through weight nine by rational
row reduction in standard Lyndon brackets. It verifies antisymmetry,
$J=0$, the averaged zero-sum equivariant conjugators, and both
projections of the two lifts for every basis vector. It also applies
the construction to the specified integral vector $q_7$, with its
actual coefficients from B.12. Rational associative-word certificates
for all these conjugators are saved in \texttt{b\_relative47.json};
the dimension table and executed output are saved as CSV and
\texttt{.out}. This exhausts the indicated rational rank-two spaces,
not merely selected vectors. The much larger rank-seven dimensions
come from the exact character identities just proved; no expanded
rank-seven matrix or sampled modular rank is claimed. The construction
in B.26 supplies the source tangentiality identities in all weights.
\end{proof}

\begin{statement}{B.28}{Open}{The prescribed profinite word remains a separate test}
The preceding results do not determine whether a prescribed profinite
loop automorphism satisfies the normalized \S47 equation
\[
 \xi=\bigl(\varphi(\xi)\sigma(\varphi(\xi))^{-1}\bigr)x^\nu,
 \qquad \sigma=\operatorname{Int}(s),\quad \nu=(\mu-1)/2.
\]
They do not construct a full profinite lift with leading term $q_7$
or decide this equation for such a lift if one is supplied. They also
do not establish that all pentagon solutions satisfy it.
These are the residual questions of LM2 B.m3 and E.3, anchored at
\S47, folios 416--420, scans 231--240. Their literal profinite
word evaluation has not been computed or reconstructed from the
associated graded data here.
\end{statement}
The lifts in B.12 are pro-$\ell$ lifts for $\ell\ge5$. They do not by
themselves give a full profinite word on which this $\varphi$ can be
evaluated: the quotient $v:F_0\to S_3$ does not factor through the
absolute pro-$\ell$ group for these primes. The $q_7$ certificate in
B.27 concerns its leading rational derivation.
The precise distinction is between a specified action on the source
relative completion and the existence of some peripheral graded action
on its radical. In a profinite solution the action on the open subgroup
$N$ is induced by the given action on $F_0$; it is not freely chosen
after passage to $U(F_0)$. B.26 deliberately allows that extra choice
and shows that the resulting existential condition is automatic. It
therefore cannot supply the missing implication for the prescribed
action. The absolute-completion obstruction of LM2 E.4 has been
replaced by the correctly typed rank-seven diagram, and that diagram
has been solved at the associated-graded level. The remaining work is
to specify and compute the actual relative action induced by a full
profinite word, or prove a comparison that recovers it from additional
data. Simultaneous filtered peripheral formality and passage from a
graded relative lift to such a profinite action have not been assumed.

